\documentclass[10pt,landscape]{article}

% Compilar com LuaLaTeX (ou XeLaTeX): lualatex naoparametrica.tex
\usepackage{fontspec}
\IfFontExistsTF{Times New Roman}
  {\setmainfont{Times New Roman}}
  {\setmainfont{TeX Gyre Termes}}

\usepackage{multicol}
\usepackage{calc}
\usepackage{ifthen}
\usepackage[landscape]{geometry}
\usepackage{graphicx}
\usepackage{amsmath, amssymb, amsthm}
\usepackage{array}
\usepackage{booktabs}
\usepackage{enumitem}
\setlist[description]{leftmargin=0pt,itemsep=3pt,parsep=0pt,topsep=0pt}
\setlist[itemize]{leftmargin=8pt,itemsep=0pt,parsep=0pt,topsep=1pt}
\usepackage[pdfauthor={Guilherme Marthe},
            pdftitle={Formulário - Estatística Não-Paramétrica},
            pdfsubject={Aulas 1 a 12},
            pdfkeywords={nonparametric} {cheatsheet}]{hyperref}
\usepackage[open,openlevel=2]{bookmark}

% Notação
\newcommand{\E}{\mathbb{E}}
\newcommand{\V}{\mathbb{V}}
\renewcommand{\P}{\mathbb{P}}
\newcommand{\I}{\mathbb{I}}
\newcommand{\R}{\mathbb{R}}
\newcommand{\N}{\mathcal{N}}
\newcommand{\cov}{\textrm{Cov}}
\newcommand{\IF}{\textrm{IF}}
\newcommand{\VC}{\textrm{VC}}
\newcommand{\MSE}{\textrm{MSE}}
\newcommand{\MISE}{\textrm{MISE}}
\newcommand{\CV}{\textrm{CV}}
\newcommand{\tr}{\textrm{tr}}
\newcommand{\Fn}{\hat{F}_n}
\newcommand{\X}{\mathbf{X}}
\newcommand{\qc}{\xrightarrow{q.c.}}
\newcommand{\dto}{\xrightarrow{d}}
\newcommand{\pto}{\xrightarrow{P}}

% Geometria e títulos (como em inferencia.tex)
\geometry{top=1cm,left=1cm,right=1cm,bottom=1cm}
\pagestyle{empty}
\makeatletter
\renewcommand{\section}{\@startsection{section}{1}{0mm}%
                                {-1ex plus -.5ex minus -.2ex}%
                                {0.5ex plus .2ex}%
                                {\normalfont\small\bfseries}}
\renewcommand{\subsection}{\@startsection{subsection}{2}{0mm}%
                                {-1explus -.5ex minus -.2ex}%
                                {0.5ex plus .2ex}%
                                {\normalfont\footnotesize\bfseries}}
\renewcommand{\subsubsection}{\@startsection{subsubsection}{3}{0mm}%
                                {-1ex plus -.5ex minus -.2ex}%
                                {1ex plus .2ex}%
                                {\normalfont\scriptsize\bfseries}}
\makeatother
\setcounter{secnumdepth}{0}
\setlength{\parindent}{0pt}
\setlength{\parskip}{0pt plus 0.5ex}
\newcommand{\secrule}{\smallskip \hrule height 2pt \smallskip}

\begin{document}
\raggedright
\scriptsize
\begin{multicols}{3}
\setlength{\premulticols}{1pt}
\setlength{\postmulticols}{1pt}
\setlength{\multicolsep}{1pt}
\setlength{\columnsep}{2pt}

{\normalsize\bfseries Formulário -- Estatística Não-Paramétrica (Aulas 1--12)}\par
\smallskip

%=====================================================================
\section{Ferramentas}\secrule
\begin{description}
  \item[Markov / Chebyshev] $\P(|T-\theta|\ge t)\le \E|T-\theta|^r/t^r$, $r>0$. Com $r=2$ e $\E T=\theta$: $\P(|T-\theta|\ge t)\le \V(T)/t^2$.
  \item[Hoeffding] $Y_i$ indep., $Y_i\in[a,b]$: $\P(|\bar Y-\E\bar Y|\ge\varepsilon)\le 2\exp\!\left(-2n\varepsilon^2/(b-a)^2\right)$. Para sinais $\sigma_i a_i\in[-1,1]$: $2e^{-n\varepsilon^2/2}$.
  \item[Borel-Cantelli I] $\sum_n\P(A_n)<\infty\Rightarrow \P(A_n\text{ i.v.})=0$. \textbf{Uso:} se $\sum_n \P(|T_n-\theta|>\varepsilon)<\infty$ para todo $\varepsilon>0$, então $T_n\qc\theta$. Toda cota exponencial em $n$ vira convergência q.c.
  \item[Esperança via cauda] $Z\ge0\Rightarrow \E Z=\int_0^\infty \P(Z\ge t)\,dt$.
  \item[LGN forte / TCL] $\E|X_1|<\infty\Rightarrow\bar X\qc\mu$; $\sigma^2<\infty\Rightarrow\sqrt n(\bar X-\mu)\dto N(0,\sigma^2)$.
  \item[Slutsky] $X_n\dto X$, $Y_n\pto c\Rightarrow X_n+Y_n\dto X+c$, $Y_nX_n\dto cX$, $X_n/Y_n\dto X/c$ ($c\ne0$).
  \item[Método delta] $\sqrt n(Y_n-\mu)\dto N(0,\Sigma)$, $g$ diferenciável em $\mu$ $\Rightarrow \sqrt n(g(Y_n)-g(\mu))\dto N(0,\nabla g^T\Sigma\nabla g)$. Ex.: $\sqrt{\Fn(x)}$: variância assintótica $\frac{F(1-F)}{4F}=\frac{1-F(x)}{4}$.
  \item[Decomposição do risco] $\MSE(\hat\theta)=(\E\hat\theta-\theta)^2+\V(\hat\theta)$.
  \item[Cauchy-Schwarz] $\left(\sum a_ib_i\right)^2\le\sum a_i^2\sum b_i^2$; $\ \sum\ell_i^2\le \max_i|\ell_i|\sum_i|\ell_i|$.
\end{description}

\subsection{Movimentos para integrais (viés/variância)}
\begin{itemize}
  \item Escrever $\E$ como integral: $\E[h^{-1}K((X-x)/h)]=\int h^{-1}K((y-x)/h)f(y)\,dy$.
  \item Mudança de variáveis $u=(y-x)/h$, $dy=h\,du$: $=\int K(u)f(x+uh)\,du$.
  \item Centrar com $\int K=1$: $\int K(u)f(x+uh)du-f(x)=\int K(u)[f(x+uh)-f(x)]du$.
  \item Módulo para dentro: $|\int g|\le\int|g|$; cotar o integrando com a suavidade.
  \item Fubini (integrando $\ge0$ ou integrável) para fatorar em kernels-produto.
  \item $\V(Z)\le\E Z^2$; $\V(\bar Z)=\V(Z_1)/n$ para i.i.d.
\end{itemize}

\subsection{Banda ótima (forma geral)}
$g(h)=Ah^{2\beta}+\dfrac{B}{nh^{d}}$ é convexa e minimizada em
$h^*=\left(\dfrac{dB}{2\beta A\,n}\right)^{\frac1{2\beta+d}}$, com $g(h^*)\asymp n^{-\frac{2\beta}{2\beta+d}}$. \\
Caso $\beta=d=1$ (Aula 1, Lema 4): $g=L^2h^2+C/(nh)$, $h_C=(C/(2L^2n))^{1/3}$, $g(h_C)=3\cdot2^{-2/3}(LC)^{2/3}n^{-2/3}$. \\
\textbf{Equilíbrio rápido:} igualar $h^{2\beta}=1/(nh^d)\Rightarrow h\asymp n^{-1/(2\beta+d)}$.

%=====================================================================
\section{Distribuição empírica (Aula 2)}\secrule
\begin{description}
  \item[Definição] $\Fn(x)=n^{-1}\sum_i\I(X_i\le x)$; $\ n\Fn(x)\sim\textrm{Bin}(n,F(x))$.
  \item[Momentos] $\E\Fn(x)=F(x)$, $\V\Fn(x)=\frac{F(x)(1-F(x))}{n}\le\frac1{4n}$; $\ \cov(\Fn(x),\Fn(y))=\frac{F(x\wedge y)-F(x)F(y)}{n}$.
  \item[Pontual] $\Fn(x)\qc F(x)$ e $\Fn(x^-)\qc F(x^-)$ (LGN).
  \item[Erro máximo] $D_n=\sup_x|\Fn(x)-F(x)|$ (distância de Kolmogorov-Smirnov).
  \item[Discretização] $x_j=\inf\{x:F(x)\ge j/k\}$, $x_0=-\infty$, $x_k=\infty$: $D_n\le\Delta_{n,k}+k^{-1}$, com $\Delta_{n,k}$ o máximo dos erros em $x_j$ e $x_j^-$ ($2(k+1)$ pontos).
  \item[Glivenko-Cantelli] $D_n\qc0$.
  \item[Velocidade] $\P(D_n\ge\varepsilon)\le\frac{16}{\varepsilon}e^{-n\varepsilon^2/2}$, $\varepsilon\in(0,1]$ $\Rightarrow D_n\le2\sqrt{n^{-1}\log n}$ q.c. para $n$ grande.
  \item[DKW (Massart)] $\P(D_n\ge\varepsilon)\le2e^{-2n\varepsilon^2}$ (constante $2$ ótima). $\E D_n\le\sqrt{\pi/(2n)}$.
  \item[Banda de confiança] $\varepsilon_n=\sqrt{\log(2/\alpha)/(2n)}$: $\P\big(\Fn-\varepsilon_n\le F\le\Fn+\varepsilon_n\ \forall x\big)\ge1-\alpha$ (truncar em $[0,1]$).
  \item[Classe GC] $\mathcal G$ é GC para $P$ se $\|P_n-P\|_{\mathcal G}=\sup_{g\in\mathcal G}|P_ng-Pg|\qc0$; $P_ng=n^{-1}\sum g(X_i)$.
\end{description}

\subsection{Moldes de consistência}
\begin{description}
  \item[Funcional Lipschitz] $|T(G_1)-T(G_2)|\le L\|G_1-G_2\|_\infty\Rightarrow|T(\Fn)-T(F)|\le LD_n\qc0$, com taxa $O(\sqrt{\log n/n})$ q.c.
  \item[Exemplos] $T(F)=\int x\,dF$ com $|X|\le M$: $\int x\,dF=\int_0^M(1-F)-\int_{-M}^0F$, logo $L=2M$. $T(F)=\int_0^MF^2$: $|G_1^2-G_2^2|=|G_1+G_2||G_1-G_2|\le2|G_1-G_2|\Rightarrow L=2M$.
  \item[Inversão do quantil] Se $f\ge C_1$ entre $x$ e $y$: $|F(x)-F(y)|\ge C_1|x-y|$. Mediana amostral ($\Fn(\hat m^-)\le\frac12\le\Fn(\hat m)$): $|\Fn(\hat m)-\tfrac12|\le n^{-1}$ (sem empates) e
  $C_1|\hat m_n-m|\le|F(\hat m)-\Fn(\hat m)|+|\Fn(\hat m)-F(m)|\le D_n+n^{-1}$.
\end{description}

%=====================================================================
\section{Dimensão VC (Aula 3)}\secrule
\begin{description}
  \item[Medida empírica] $P_n(A)=n^{-1}\sum\I(X_i\in A)$; $\mathcal C$ é GC se $\|P_n-P\|_{\mathcal C}\qc0$. Contraexemplo: todos os borelianos com $P$ contínua ($C=\{X_1..X_n\}$).
  \item[Traço] $\mathcal C_A=\{C\cap A:C\in\mathcal C\}$; $\mathcal C$ \emph{fragmenta} $A$ se $\mathcal C_A=2^A$.
  \item[Coef. de fragmentação] $s(\mathcal C,n)=\max_{|A|=n}|\mathcal C_A|\le2^n$.
  \item[Dimensão VC] $\VC(\mathcal C)=\sup\{n:s(\mathcal C,n)=2^n\}$. Para $\VC=d$: (i) exibir \emph{um} conjunto de $d$ pontos fragmentado; (ii) mostrar que \emph{nenhum} conjunto de $d+1$ pontos é fragmentado.
  \item[Sauer-Shelah] $\VC=d<\infty\Rightarrow s(\mathcal C,n)\le\sum_{i=0}^d\binom ni\le(n+1)^d$.
  \item[Desig. de VC] $\P(\|P_n-P\|_{\mathcal C}>\varepsilon)\le8\,s(\mathcal C,n)\,e^{-n\varepsilon^2/32}$.
  \item[Prova (cadeia)] Simetrização ($n\varepsilon^2\ge2$): $\le2\P(\|P_n-P_n'\|>\frac\varepsilon2)$ $\to$ sinais: $\le4\P(\sup|\frac1n\sum\sigma_i\I(X_i\in C)|>\frac\varepsilon4)$ $\to$ contagem $+$ Hoeffding: $\le8s(\mathcal C,n)e^{-n\varepsilon^2/32}$.
  \item[Corolário] Classe VC $\Rightarrow$ GC universal (polinomial $\times$ exponencial é somável; Borel-Cantelli). Semirretas ($\VC=1$) $\Rightarrow$ Glivenko-Cantelli.
  \item[Radon] $d+2$ pontos em $\R^d$ admitem partição $I,J$ com $\textrm{conv}(x_I)\cap\textrm{conv}(x_J)\ne\emptyset$ (cota superior para conjuntos convexos/semiespaços).
  \item[Operações] $\{A\cup B\}$ ou $\{A\cap B\}$ com $A\in\mathcal C_1,B\in\mathcal C_2$: $s\le s(\mathcal C_1,n)\,s(\mathcal C_2,n)$; complementos preservam $\VC$; $\mathcal C\subseteq\mathcal C'\Rightarrow\VC(\mathcal C)\le\VC(\mathcal C')$.
\end{description}
\begin{center}
\begin{tabular}{lc}
\toprule
Classe & VC \\
\midrule
Semirretas $(-\infty,x]$ em $\R$ & $1$ \\
Intervalos $[a,b]$ em $\R$ & $2$ \\
Retângulos (lados paralelos) em $\R^2$ & $4$ \\
Semiespaços $\{\langle w,x\rangle\ge b\}$ em $\R^d$ & $d+1$ \\
Discos em $\R^2$ & $3$ \\
\bottomrule
\end{tabular}
\end{center}
\textbf{Técnica em $\R$:} ordene $x_1<\cdots<x_k$ e liste os padrões possíveis do traço; ache um rótulo impossível (ex.: incluir $x_2,x_4$ e excluir $x_1,x_3$).

%=====================================================================
\section{Funcionais estatísticos (Aula 4)}\secrule
\begin{description}
  \item[Funcional / plug-in] $T:\mathcal D\to\R$; plug-in $T(\Fn)$.
  \item[Linear] $T(F)=\int g\,dF\Rightarrow T(\Fn)=n^{-1}\sum g(X_i)$, não-viesado, $\V=\V[g(X_1)]/n$, $\sqrt n(T(\Fn)-T(F))\dto N(0,\V g(X_1))$.
  \item[Composição] $T=h(\int g_1dF,\dots,\int g_kdF)$, $h$ contínua $\Rightarrow T(\Fn)\qc T(F)$. Ex.: variância, $h(a,b)=b-a^2$. Quantis \emph{não} são composições.
  \item[Plug-in da variância] $\hat\sigma^2=n^{-1}\sum(X_i-\bar X)^2$, $\E\hat\sigma^2=\frac{n-1}n\sigma^2$.
  \item[Contaminação] $F_{\varepsilon,x}=(1-\varepsilon)F+\varepsilon\delta_x$.
  \item[Função de influência] $\IF(x)=\lim_{\varepsilon\downarrow0}\frac{T(F_{\varepsilon,x})-T(F)}{\varepsilon}=\frac{d}{d\varepsilon}T(F_{\varepsilon,x})\big|_{\varepsilon=0}$. Sempre $\int\IF\,dF=0$ (checagem!).
\end{description}

\begin{center}
\begin{tabular}{ll}
\toprule
$T(F)$ & $\IF(x)$ \\
\midrule
$\int g\,dF$ & $g(x)-T(F)$ \\
$\mu$ & $x-\mu$ \\
$F(b)-F(a)$ & $\I(a<x\le b)-T(F)$ \\
$\sigma^2$ & $(x-\mu)^2-\sigma^2$ \\
$\xi_p$ (quantil) & $\dfrac{p-\I(x\le\xi_p)}{f(\xi_p)}$ \\
mediana & $\dfrac{\textrm{sinal}(x-m)}{2f(m)}$ \\
\bottomrule
\end{tabular}
\end{center}

\subsection{Cálculo de IF}
\begin{description}
  \item[Regra da cadeia] $T=h(\theta_1,..,\theta_k)$, $\theta_j=\int g_jdF$: $\IF(x)=\sum_j\partial_jh(\theta)\,(g_j(x)-\theta_j)$.
  \item[Produto/quociente/potência] $\IF_{T_1T_2}=T_2\IF_1+T_1\IF_2$; $\ \IF_{T_1/T_2}=\frac{\IF_1-(T_1/T_2)\IF_2}{T_2}$; $\ \IF_{T_1^a}=aT_1^{a-1}\IF_1$.
  \item[Momento centrado] $\mu_k(F)=\int(y-\mu(F))^kdF$: derivar a integral \emph{e} o centro:
  $\IF_{\mu_k}(x)=(x-\mu)^k-\mu_k-k\mu_{k-1}(x-\mu)$. ($k=2$: $\mu_1=0$; $k=3$: $-3\sigma^2(x-\mu)$.)
  \item[Assimetria] $T=\mu_3/\sigma^3$: $\IF=\dfrac{(x-\mu)^3-\mu_3-3\sigma^2(x-\mu)}{\sigma^3}-\dfrac{3T}{2}\cdot\dfrac{(x-\mu)^2-\sigma^2}{\sigma^2}$. Requer $\int x^6dF<\infty$ para $\tau^2<\infty$.
  \item[Quantil (truque)] derivar $(1-\varepsilon)F(\xi_p(\varepsilon))+\varepsilon\I(x\le\xi_p(\varepsilon))=p$ em $\varepsilon=0$.
\end{description}

\subsection{Normalidade assintótica}
\begin{description}
  \item[(A1) von Mises] $T(G)-T(F)=\int\IF\,d(G-F)+R(G,F)$, $\sqrt nR(\Fn,F)\pto0$, $\tau^2=\int\IF^2dF\in(0,\infty)$.
  \item[Teorema] $\sqrt n(T(\Fn)-T(F))=\frac1{\sqrt n}\sum\IF(X_i)+o_P(1)\dto N(0,\tau^2)$.
  \item[EP e IC] $\hat\tau^2=n^{-1}\sum\widehat{\IF}(X_i)^2$; $\ T(\Fn)\pm z_{1-\alpha/2}\hat\tau/\sqrt n$.
  \item[Mediana] $\sqrt n(\hat m-m)\dto N\!\left(0,\frac1{4f(m)^2}\right)$ (exige estimar $f$).
  \item[Robustez] sensibilidade a erros grosseiros $\sup_x|\IF(x)|$: média $\infty$, mediana $1/(2f(m))$.
\end{description}

%=====================================================================
\section{Jackknife e Bootstrap (Aula 6)}\secrule
\begin{description}
  \item[Exclusão] $\hat\theta_{-i}=T(\hat F_{n,-i})$, $\bar\theta=n^{-1}\sum\hat\theta_{-i}$.
  \item[Jackknife] $\hat b_{\textrm{jack}}=(n-1)(\bar\theta-\hat\theta_n)$, $\ \hat\theta_{\textrm{jack}}=n\hat\theta_n-(n-1)\bar\theta$.
  \item[(A1)] $\E\hat\theta_n=\theta+\frac an+\frac b{n^2}+O(n^{-3})$ $\Rightarrow$ $\E\hat b_{\textrm{jack}}=\frac an+O(n^{-2})$, $\E\hat\theta_{\textrm{jack}}-\theta=O(n^{-2})$. Se o viés é \emph{exatamente} $a/n$: $\hat\theta_{\textrm{jack}}$ é não-viesado.
  \item[Pseudo-valores] $\tilde\theta_i=n\hat\theta_n-(n-1)\hat\theta_{-i}$; $\hat\theta_{\textrm{jack}}=\overline{\tilde\theta}$; $\ \hat v_{\textrm{jack}}=\frac{n-1}n\sum(\hat\theta_{-i}-\bar\theta)^2=\frac1{n(n-1)}\sum(\tilde\theta_i-\overline{\tilde\theta})^2$.
  \item[Identidades] $\bar X_{-i}=\frac{n\bar X-X_i}{n-1}$; $\ \bar X_{-i}-\bar X=\frac{\bar X-X_i}{n-1}$; $\ \sum_i(\bar X_{-i}-\bar X)^2=\frac{n\hat\sigma^2}{(n-1)^2}$; $\ \E\bar X^2=\mu^2+\frac{\sigma^2}n$.
  \item[Média] $\hat b_{\textrm{jack}}=0$, $\hat v_{\textrm{jack}}=S^2/n$ ($S^2$ com $n-1$).
  \item[$\mu^2$] $\hat\theta=\bar X^2$: $a=\sigma^2,b=0$; $\bar\theta=\frac1n\sum\bar X_{-i}^2=\bar X^2+\frac{\hat\sigma^2}{(n-1)^2}$, $\hat\theta_{\textrm{jack}}=\bar X^2-\frac{\hat\sigma^2}{n-1}$ (não-viesado).
\end{description}
\begin{description}
  \item[Amostra bootstrap] $X_1^*,\dots,X_n^*$ i.i.d. $\sim\Fn$; $\hat\theta^*=T(\Fn^*)$. O bootstrap é o plug-in de funcionais como $\textrm{se}_F(\hat\theta)$.
  \item[Viés / EP] $\hat b^*=\E^*\hat\theta^*-\hat\theta_n$; corrigido $2\hat\theta_n-\bar\theta^*$. $\hat{\textrm{se}}^2_{\textrm{boot}}=\frac1B\sum_b(\hat\theta^*_b-\bar\theta^*)^2$ (Monte Carlo para $\V^*$).
  \item[Momentos sob $\Fn$] $\E^*\bar X^*=\bar X$, $\V^*\bar X^*=\hat\sigma^2/n$, $\E^*(\bar X^*)^2=\bar X^2+\hat\sigma^2/n$. Para $\bar X^2$: $\hat b^*=\hat\sigma^2/n$; $\bar X^2-\hat\sigma^2/n$ tem viés $\sigma^2/n^2$.
  \item[Intervalo t] $\hat\theta_n\pm z_{1-\alpha/2}\hat{\textrm{se}}_{\textrm{boot}}$.
  \item[Percentil] $[q^*_{\alpha/2},q^*_{1-\alpha/2}]$ (quantis das réplicas).
  \item[Consistência] $\V X_1<\infty$, média: $\sup_x|\P^*(\sqrt n(\hat\theta^*-\hat\theta_n)\le x)-\P(\sqrt n(\hat\theta_n-\theta)\le x)|\qc0$.
  \item[Contagem] amostras ordenadas: $n^n$; multiconjuntos distintos: $\binom{2n-1}{n}=\frac{(2n-1)!}{n!(n-1)!}$ (estrelas e barras: $n$ bolas em $n$ caixas).
\end{description}

%=====================================================================
\section{Estimação de densidade}\secrule
\subsection{Histograma (Aula 1)}
\begin{description}
  \item[Definição] $[0,1]$ em $m$ intervalos de comprimento $h=1/m$: $\hat f_n(x)=\frac1{nh}\sum_i\I(X_i\in B_j)$, $x\in B_j$; $p_j=\int_{B_j}f$.
  \item[Momentos] $\E\hat f_n(x)=p_j/h$; $\ \V\hat f_n(x)=\frac{p_j(1-p_j)}{nh^2}$.
  \item[(A1)] $|f'|\le L$ $\Rightarrow$ $|f(u)-f(x)|\le L|u-x|$.
  \item[Viés] $|\E\hat f_n(x)-f(x)|\le Lh$. \textbf{Cota} $f\le1+L/2$ (integrar $f(x)\le f(y)+L|x-y|$ em $y$).
  \item[Variância] $\V\hat f_n(x)\le\frac{f(x)+Lh}{nh}\le\frac{1+L/2+Lh}{nh}$.
  \item[Risco] $\MSE\le L^2h^2+\frac{1+L/2}{nh}+\frac Ln$; $\ \MISE\le L^2h^2+\frac1{nh}+\frac Ln$; $h\asymp n^{-1/3}$, taxa $n^{-2/3}$.
\end{description}

\subsection{KDE (Aula 7)}
\begin{description}
  \item[Hölder $\Sigma(\beta,L)$] $\ell=$ maior inteiro $<\beta$; $f$ $\ell$ vezes diferenciável, $|f^{(\ell)}(y)-f^{(\ell)}(x)|\le L|y-x|^{\beta-\ell}$.
  \item[Taylor] $f(x+t)=\sum_{j=0}^{\ell}\frac{f^{(j)}(x)}{j!}t^j+R_x(t)$, $|R_x(t)|\le\frac L{\ell!}|t|^\beta$.
  \item[Kernel] $\int K=1$, $R(K)=\int K^2<\infty$; $\kappa_j=\int|u|^j|K(u)|du$.
  \item[Ordem $k$] $\kappa_k<\infty$ e $\int u^jK=0$, $1\le j\le k$. Simétrico $+\kappa_1<\infty\Rightarrow$ ordem $1$. Para $\Sigma(\beta,L)$ usar ordem $\ell$; $\beta\le1$: nenhum momento precisa anular.
  \item[KDE] $\hat f_n(x)=\frac1{nh}\sum_iK\!\left(\frac{X_i-x}h\right)$.
  \item[Momentos] $\E\hat f_n(x)=\int K(u)f(x+uh)du$; $\ \V\hat f_n(x)\le\frac1{nh}\int K(u)^2f(x+uh)du\le\frac{\sup f\,R(K)}{nh}$.
  \item[Viés] $\E\hat f_n(x)-f(x)=\int K(u)R_x(uh)du$, $|\cdot|\le\frac{L\kappa_\beta}{\ell!}h^\beta$.
  \item[$f$ limitada] $K$ limitado de ordem $\ell$, $\kappa_\beta<\infty$: $f(x)=\int K(u)f(x+u)du-\int KR_x$ $\Rightarrow f\le\sup|K|+\frac{L\kappa_\beta}{\ell!}$.
  \item[Risco] $\MSE(\hat f_n(x))\le C\left(h^{2\beta}+\frac1{nh}\right)$; $h=n^{-1/(2\beta+1)}\Rightarrow\sup_x\MSE\le2Cn^{-\frac{2\beta}{2\beta+1}}$. ($\beta=1$: $n^{-2/3}$; $\beta=2$: $n^{-4/5}$.)
\end{description}
\begin{center}
\begin{tabular}{lccc}
\toprule
Kernel & $K(u)$ & $\kappa_2$ & $R(K)$ \\
\midrule
Uniforme & $\frac12\I(|u|\le1)$ & $\frac13$ & $\frac12$ \\
Epanechnikov & $\frac34(1-u^2)\I(|u|\le1)$ & $\frac15$ & $\frac35$ \\
Gaussiano & $(2\pi)^{-1/2}e^{-u^2/2}$ & $1$ & $\frac1{2\sqrt\pi}$ \\
\bottomrule
\end{tabular}
\end{center}

\subsection{Regra geral de taxas}
\begin{description}
  \item[$d$ dimensões, derivada $s$] viés $\asymp h^{\beta-s}$, variância $\asymp(nh^{d+2s})^{-1}$ $\Rightarrow h^*\asymp n^{-\frac1{2\beta+d}}$, $\MSE\asymp n^{-\frac{2(\beta-s)}{2\beta+d}}$.
  \item[KDE bivariado] $\hat f_h(x,z)=\frac1{nh^2}\sum K(\frac{X_i-x}h)K(\frac{Z_i-z}h)$; $\E\hat f_h=\iint K(u)K(v)f(x+uh,z+vh)\,du\,dv$.
  Com $|f(x,z)-f(x',z')|\le L(|x-x'|^\beta+|z-z'|^\beta)$, $\beta\le1$: viés $\le2L\kappa_0\kappa_\beta h^\beta$; variância $\le\frac{\sup f\,R(K)^2}{nh^2}$; $f\le1+\frac{2L}{\beta+1}$ (integrar em $[x,x+1]\times[z,z+1]$). $h\asymp n^{-\frac1{2\beta+2}}$, $\MSE\asymp n^{-\frac{\beta}{\beta+1}}$ (pior que $n^{-\frac{2\beta}{2\beta+1}}$: maldição da dimensão).
  \item[Derivada $s$ (Tsybakov 1.2)] $\hat p_{n,s}=\frac1{nh^{s+1}}\sum K(\frac{X_i-x}h)$, $\int u^sK=s!$: viés $\le ch^{\beta-s}$, $\V\le c'(nh^{2s+1})^{-1}$.
\end{description}

\subsection{Estimação não-viesada do risco (Aula 8)}
\begin{description}
  \item[Decomposição] $\MISE(h)=A(h)-2B(h)+\int f^2$, $A=\E\int\hat f_h^2$, $B=\E\int\hat f_hf$. Minimizar $J(h)=A-2B$.
  \item[Termos] $\int\hat f_h^2$ é não-viesado para $A$. $\ B(h)=\E\hat f_h(X_{n+1})=\frac1h\E K\!\left(\frac{X_2-X_1}h\right)$.
  \item[Estimador ingênuo] $\E\left[\frac1n\sum\hat f_h(X_i)\right]=\frac{K(0)}{nh}+\frac{n-1}nB(h)$ (viesado).
  \item[LOO] $\hat f_{h,-i}(x)=\frac1{(n-1)h}\sum_{j\ne i}K(\frac{X_j-x}h)$.
  \item[CV] $\CV(h)=\int\hat f_h^2-\frac2n\sum_i\hat f_{h,-i}(X_i)$; $\ \E\CV(h)=J(h)=\MISE(h)-\int f^2$.
  \item[Cálculo] $\bar K(v)=\int K(t)K(t+v)dt$: $\CV(h)=\frac1{n^2h}\sum_{i,j}\bar K(\frac{X_j-X_i}h)-\frac2{n(n-1)h}\sum_{i\ne j}K(\frac{X_j-X_i}h)$. Gaussiano: $\bar K=$ densidade $N(0,2)$.
\end{description}

%=====================================================================
\section{Regressão: linear smoothers (Aulas 9--10)}\secrule
\begin{description}
  \item[Modelo] $Y_i=r(X_i)+\varepsilon_i$, $r(x)=\E[Y\mid X=x]$, $\E[\varepsilon\mid X]=0$, $\V[\varepsilon\mid X]=\sigma^2$.
  \item[MQ] $\hat\beta=(\mathbb X^T\mathbb X)^{-1}\mathbb X^T\mathbf Y$; $\ell_i(x)=x^T(\mathbb X^T\mathbb X)^{-1}X_i$, $\sum\ell_i=1$ (com intercepto).
  \item[Ridge] $\hat\beta_\lambda=(\mathbb X^T\mathbb X+\lambda I)^{-1}\mathbb X^T\mathbf Y$.
  \item[Linear smoother] $\hat r_n(x)=\sum_i\ell_i(x)Y_i$, $\ell$ depende de $x$ e dos $X_i$, não dos $Y_i$.
  \item[Matriz] $\mathbf L_{ij}=\ell_j(X_i)$, $\hat{\mathbf r}=\mathbf L\mathbf Y$; $\nu=\tr(\mathbf L)$. MQ: $\nu=p$.
\end{description}
\begin{center}
\begin{tabular}{ll}
\toprule
Método & $\ell_i(x)$ \\
\midrule
Regressograma & $\I(X_i\in B_j)/n_j$, $x\in B_j$ \\
$k$-NN & $\I(i\in N_k(x))/k$ \\
Médias locais & $\I(|X_i-x|\text{ em }(-h,h])/\sum_s(\cdot)$ \\
Nadaraya-Watson & $K(\frac{X_i-x}h)/\sum_sK(\frac{X_s-x}h)$ \\
Pol. local grau $q$ & $\frac1{nh}K(Z_i)e_1^TB_x^{-1}U(Z_i)$ \\
Proc. gaussiano & $[\mathbf c(x)^T(\mathbf C+\sigma^2I)^{-1}]_i$ \\
\bottomrule
\end{tabular}
\end{center}
\begin{description}
  \item[Pol. local] $Z_i=\frac{X_i-x}h$, $U(u)=(1,u,\frac{u^2}{2!},..,\frac{u^q}{q!})^T$, $B_x=\frac1{nh}\sum K(Z_i)U(Z_i)U(Z_i)^T$. Reproduz polinômios: $\sum\ell_i=1$, $\sum\ell_i(X_i-x)^j=0$, $1\le j\le q$. $q=0$: Nadaraya-Watson.
  \item[Cond. ao desenho] $\E[\hat r_n(x)\mid\X]=\sum\ell_i(x)r(X_i)$; $\ \V[\hat r_n(x)\mid\X]=\sigma^2\|\ell(x)\|^2$. Se $\sum\ell_i=1$: viés $=\sum\ell_i(x)(r(X_i)-r(x))$.
\end{description}

\subsection{Risco de linear smoothers}
\begin{description}
  \item[(A1)] $r\in\Sigma(\beta,L)$, $m=$ maior inteiro $<\beta$.
  \item[(A2)] (P1) $\sum\ell_i=1$, $\sum\ell_i(X_i-x)^j=0$, $1\le j\le m$; (P2) $\ell_i(x)=0$ se $|X_i-x|>h$; (P3) $\sum|\ell_i|\le C_0$, $\max|\ell_i|\le\frac{C_0}{nh}$.
  \item[Teorema] viés $\le\frac{LC_0}{m!}h^\beta$, $\ \V\le\frac{\sigma^2C_0^2}{nh}$, risco cond. $\le C(h^{2\beta}+\frac1{nh})$; $h=n^{-1/(2\beta+1)}$: $2Cn^{-\frac{2\beta}{2\beta+1}}$. (P1)$\leftrightarrow$ordem do kernel; (P2)$\leftrightarrow$suporte; (P3)$\leftrightarrow\kappa_\beta,R(K)$.
  \item[Regressograma] $r$ $L$-Lipschitz, $n_j\ge c_0nh$: viés $\le Lh$, $\V\le\frac{\sigma^2}{n_j}\le\frac{\sigma^2}{c_0nh}$, $\MSE\to0$ se $h_n\to0$ e $nh_n\to\infty$.
  \item[Nadaraya-Watson] $0\le K\le K_{\max}$, suporte $[-1,1]$, $\beta\le1$, $\sum_sK(\frac{X_s-x}h)\ge c_0nh$: $C_0=\max(1,K_{\max}/c_0)$.
  \item[Pol. local] $q\ge m$; (D1) $\lambda_{\min}(B_x)\ge\lambda_0$; (D2) $\#\{|X_i-x|\le h\}\le a_0nh$: $|\ell_i|\le\frac{K_{\max}\sqrt{q+1}}{\lambda_0nh}$. $B_x\to f(x)\mathbf B$, $\mathbf B=\int KUU^T$.
  \item[Kernel equivalente] $K^*(u)=e_1^T\mathbf B^{-1}U(u)K(u)$ é kernel de ordem $q$.
\end{description}

\subsection{Escolha de banda e inferência}
\begin{description}
  \item[CV] $\CV=\frac1n\sum(Y_i-\hat r_{-i}(X_i))^2$.
  \item[Atalho] Pesos normalizados $\ell_j=g_j/\sum_sg_s$ ou MQ: $Y_i-\hat r_{-i}(X_i)=\frac{Y_i-\hat r_n(X_i)}{1-\mathbf L_{ii}}$.
  \item[Sherman-Morrison] $(A-vv^T)^{-1}=A^{-1}+\frac{A^{-1}vv^TA^{-1}}{1-v^TA^{-1}v}$.
  \item[GCV] $\frac1n\sum\left(\frac{Y_i-\hat r_n(X_i)}{1-\nu/n}\right)^2$.
  \item[Risco no desenho] $\mathcal R_n=\frac1n\E[\|\hat{\mathbf r}-\mathbf r\|^2\mid\X]$.
  \item[Formas quadráticas] $\E[\|\mathbf v+M\boldsymbol\varepsilon\|^2\mid\X]=\|\mathbf v\|^2+\sigma^2\tr(M^TM)$.
  \item[$C_p$ de Mallows] $\hat C_p=\frac1n\|\mathbf Y-\mathbf L\mathbf Y\|^2+\frac{2\sigma^2\nu}n$, $\E[\hat C_p\mid\X]=\mathcal R_n+\sigma^2$.
  \item[IC] erros $N(0,\sigma^2)$: $\hat r_n(x)\pm z_{1-\alpha/2}\,\sigma\|\ell(x)\|$ cobre $\E[\hat r_n(x)\mid\X]$ (não $r(x)$: viés!).
\end{description}

%=====================================================================
\section{Receitas para a prova}\secrule
\begin{description}
  \item[Consistência q.c.] Lipschitz em $\|\cdot\|_\infty$ $\Rightarrow|T(\Fn)-T(F)|\le LD_n$ $\Rightarrow$ GC/DKW. Taxa via $D_n\le2\sqrt{\log n/n}$.
  \item[VC] Fragmentar $d$ pontos explícitos; para $d+1$, ordenar e achar rótulo impossível; concluir GC por Sauer + desigualdade VC.
  \item[IF $\to$ IC] Calcular IF (cadeia/quociente), checar $\int\IF dF=0$, $\tau^2=\E\IF^2$, momentos necessários, IC de Wald.
  \item[Viés/variância/MSE] (1) $\E$ como integral; (2) $u=(y-x)/h$; (3) centrar com $\int K=1$; (4) Taylor/Hölder; (5) variância $\le\E Z^2$; (6) equilibrar $h$.
  \item[Jackknife] calcular $\hat\theta_{-i}$ via $\bar X_{-i}$, somar, comparar com (A1).
  \item[Armadilhas] $\hat\sigma^2$ ($n$) vs.\ $S^2$ ($n-1$); $\ell$ vs.\ $m$ (inteiro \emph{estritamente} menor); taxa depende de $d$; resultados de regressão são \emph{condicionais} a $\X$.
\end{description}

\end{multicols}
\end{document}
